\documentclass[10pt]{article}
\usepackage[margin=0.85in,top=0.7in,bottom=0.7in]{geometry}
\usepackage[T1]{fontenc}
\usepackage{lmodern}
\usepackage{microtype}
\usepackage{xcolor}
\usepackage{hyperref}
\hypersetup{colorlinks=true,urlcolor=blue!50!black}
\usepackage{enumitem}
\setlist[itemize]{itemsep=2pt,topsep=3pt}
\setlength{\parindent}{0pt}
\setlength{\parskip}{0.6em}
\pagestyle{empty}
\begin{document}
\begin{flushright}
Quang-Vinh Dang\\ British University Vietnam\\ Hung Yen, Vietnam\\ \href{mailto:vinh.dq4@buv.edu.vn}{vinh.dq4@buv.edu.vn}\\ ORCID: 0000-0002-3877-8024\\[0.8em]
September 28, 2026
\end{flushright}

Editor-in-Chief\\
IEEE Transactions on Neural Networks and Learning Systems

Dear Editor-in-Chief,


I am pleased to submit the manuscript entitled \emph{``DisCo-CFL: Calibrated, Fair and Explainable Clustered Federated Learning via Disattenuated Class-Conditional Gradient Signatures''} for consideration as a regular paper in the IEEE Transactions on Neural Networks and Learning Systems.

Clustered federated learning (CFL) trains one model per group of clients that share a data distribution. A recent survey of the field identifies several open problems: the clustering signal is high-dimensional and changes from round to round, methods tend to group clients by label distribution rather than by concept, the number of clusters is set through hyperparameters without a stable meaning, performance degrades under quantity skew, and grouping decisions are neither private nor explained. The manuscript addresses these problems with a single algorithm and its analysis.

The main contributions are:
\begin{itemize}
\item \textbf{A new CFL algorithm.} Each client uploads, once, sketched class-conditional gradient signatures computed on two random halves of its data. Spearman's correction for attenuation, estimated from the two halves, yields a calibrated similarity that is invariant to label and quantity skew and equals one for clients that share a concept. A class-wise bottleneck linkage then determines the number of clusters from one threshold on this calibrated scale.
\item \textbf{Guarantees beyond accuracy.} We prove safety (conflicting clients are never merged), maximality and exact recovery with a sample complexity that does not depend on cluster sizes, which gives fairness towards minority groups, and a minimax lower bound that matches it up to a factor $1/\Delta$. We also prove differential privacy with an explicit cost, the fidelity of the class-level explanations, and the Bayes optimality of prior-corrected cluster models. The deterministic core of all proofs (25 theorems) is machine-checked in Lean~4.
\item \textbf{Transparency and accountability.} The server returns assignment certificates with uncertainty, an audit log, anomaly flags, class-level explanations of why clusters differ, and a test for newcomers with a novel concept.
\item \textbf{Extensive evaluation.} 618 training runs on Fashion-MNIST, MNIST, SVHN and CIFAR-10 against nine baselines, with robustness sweeps, ablations and significance tests. DisCo-CFL recovers the concept groups without knowing their number and merges on average 0.1\% of conflicting client pairs, against 8--84\% for baselines that are given the true number of clusters. The limitations, including the weaker results on CIFAR-10 with a short warm-up, are reported explicitly.
\end{itemize}

We believe the manuscript fits the scope of the Transactions, as it combines a new learning algorithm for distributed neural network training with theoretical guarantees and a thorough empirical study. The main paper is 8 pages long. Full proofs, the Lean~4 formalisation and complete result tables are provided in a separate supplementary file. The manuscript and the supplementary file are anonymised for double-anonymous review; code, raw results and the Lean project will be made public upon acceptance.

This manuscript is original, has not been published previously, and is not under consideration for publication elsewhere. The author declares no conflict of interest. Thank you for considering this submission. I look forward to hearing from you.

Sincerely,\\[0.8em]
Quang-Vinh Dang\\
British University Vietnam
\end{document}
